\subsection*{2.5 Vitali Set}

In this section, we prove the existence of a non-measurable set in the interval $[0,1]$, with respect to the usual length function. To do so, we define the ``modulo-1 addition'' operation, denoted by $x+y \bmod 1$, which takes the fractional part of $x+y$, ensuring that the results are always between 0 and 1. Geometrically, we identify the two endpoints of $[0,1]$ and regard the interval as a circle, where the mod-1 addition corresponds to a rotation.

We also define the \textit{translation} of a set $A$ by a distance $d$ as
\[
A+d \triangleq \{x+d \bmod 1 : x\in A\}.
\]

We will assume a technical property, namely that the measure function is translation-invariant (or rotation-invariant, if we are picturing the space as a circle). A set function $m$ is said to be \textit{translation-invariant} if $m(A)=m(A+d)$ holds for any set $A$ in the domain of $m$ and any real number $d$. We will assume that a set $A$ and any translation $A+d$ of $A$ have the same measure. This property is natural to expect because the usual arc length function on the circle is rotational-invariant.

The following theorem shows that it is impossible to define a probability measure that models the uniform distribution for all subsets of $[0,1]$. The set $V$ constructed in the proof of the theorem is known as the Vitali set.

\begin{thmbox}{2.9}
There is no set function $m:\mathcal{P}([0,1])\to [0,\infty)$ that satisfies:
\begin{enumerate}[label=\arabic*.]
    \item Monotonicity
    \item Countable additivity
    \item Translation-invariance
    \item $m([a,b])=b-a$ for all $a<b$
\end{enumerate}
\end{thmbox}

\textit{Proof} Because $0$ and $1$ are regarded as the same modulo 1, we will consider the interval $[0,1)$ in the following discussion, so that all points in $[0,1)$ are distinct modulo 1. Since we can include a single point in an arbitrarily small closed interval, by the assumptions of monotonicity, a single point has measure zero and hence is negligible.

Define a relation on $[0,1)$ by $x\sim y \Longleftrightarrow x-y\in \mathbb{Q}$. It can be shown that this is an equivalence relation, satisfying the following properties:
\begin{enumerate}[label=(\roman*)]
    \item (reflexive) $x\sim x$;
    \item (symmetric) If $x\sim y$, then $y\sim x$;
    \item (transitive) If $x\sim y$ and $y\sim z$, then $x\sim z$.
\end{enumerate}

We denote the equivalence class that contains $x$ by
\[
[x] \triangleq \{z\in [0,1) : z\sim x\}.
\]
For example, the equivalence class $[0]$ that contains $0$ consists of all rational numbers in the interval $[0,1)$.

We construct a subset of $[0,1)$ that contains exactly one representative from each equivalence class. We denote the resulting infinite set by $V$. By the Axiom of Choice, we can always choose a representative for each class, even though there are uncountably many equivalence classes. We choose $0$ to be the representative for the equivalence class $[0]=\mathbb{Q}\cap [0,1)$. Therefore, the elements of $V$ are all irrational, except for $0$.

We claim that the translations $V+r$ of $V$, with $r$ ranging over all rational numbers in $[0,1)$, form a partition of $[0,1)$, i.e.,
\begin{equation}
[0,1)=\biguplus_{r\in \mathbb{Q}\cap [0,1)} (V+r).
\tag{2.1}
\end{equation}

We first prove that the translated sets $V+r_1$ and $V+r_2$ are disjoint for any two distinct rational numbers $r_1$ and $r_2$ in $[0,1)$. Suppose for the sake of contradiction that $\alpha$ is a common element in $V+r_1$ and $V+r_2$. Then there exist $v_1$ and $v_2$ in $V$ such that $\alpha=v_1+r_1=v_2+r_2 \bmod 1$. This implies that $v_1=v_2+(r_2-r_1)\bmod 1$. Since $r_2-r_1\neq 0 \bmod 1$, this contradicts the assumption that $v_1$ and $v_2$ are the representatives of two different equivalence classes. Therefore, the sets $V+r_1$ and $V+r_2$ are disjoint.

Next, we show that any real number $x$ in $[0,1)$ belongs to $V+r$ for some rational number $r\in [0,1)$. Let $y$ be the representative of the coset that contains $x$. We distinguish two cases. If $x\ge y$, then the difference $r_0\triangleq x-y$ is a rational number in $[0,1)$, and $x\in V+r_0$ since $x$ can be written as $x=y+(x-y)=y+r_0 \bmod 1$. If $x<y$, then $r_1\triangleq 1-y+x$ is a rational number in $[0,1)$, and we have $x\in V+r_1$, as we can write $x$ as $x=y+r_1-1 \bmod 1$. This completes the proof of the claim in (2.1).

Since the rational numbers in $[0,1)$ are countable, the union in (2.1) is a countable and disjoint union. By the assumption of countable additivity and translation-invariance, we obtain
\begin{equation}
m([0,1))=m\left(\biguplus_{r\in \mathbb{Q}\cap [0,1)} (V+r)\right)
=
\sum_{r\in \mathbb{Q}\cap [0,1)} m(V).
\tag{2.2}
\end{equation}

If $m$ is defined for all sets, then the value of $m(V)$ is well-defined and equal to a constant $c$. If $c=0$, we obtain from (2.2) that $m([0,1))=0$, which contradicts the fact that $m([0,1))\ge m([0,1/2])=1/2$. On the other hand, if $c>0$, we get $m([0,1))=\infty$ from (2.2), which contradicts $m([0,1))\le m([0,1])=1$. \hfill $\square$

This theorem says that even with the simplest uniform distribution, there exists a set to which we cannot assign any probability. This means that we have to abandon some of the assumptions in Theorem 2.9. Since monotonicity, countable additivity, and translation-invariance are desirable properties, and we certainly want the probability of interval $[a,b]$ to be $b-a$ in the uniform distribution, we cannot require that a measure function is defined for all subsets. Instead, we assign probabilities only to the subsets in a $\sigma$-algebra. Subsets that are not in the $\sigma$-algebra have no probability.
